%! Setting Up --- Learning from Data — Unit 8, Lesson 8.0 — Group Activity (Answer Key)
\documentclass[10pt]{article}
\usepackage{linalg-article}
\usepackage{linalg-key}   % pulls in -boxes; do NOT also load -boxes
\newcommand{\TallMath}[1]{$\displaystyle #1\rule[-1.4em]{0pt}{3.2em}$}
\newcommand{\vv}{\mathbf{v}}
\newcommand{\bb}{\mathbf{b}}
\newcommand{\zero}{\mathbf{0}}
\newcommand{\relu}{\mathrm{ReLU}}

\begin{document}

\pageheader{Unit 8, Lesson 8.0}{Group Activity --- Answer Key}
\namepartnerperiod

\vspace{0.10in}

\begin{headlinebox}{blush}
\small\textbf{Building a Learning Function.} A neural network layer is $\relu(A\vv+\bb)$: a
\emph{linear step} (matrix times vector, plus a bias --- Unit~1) then the \emph{nonlinear step}
$\relu(x)=\max(0,x)$, applied to each entry. Stacking and shifting ReLUs builds \textbf{piecewise-linear}
functions that bend to fit data; the network \emph{learns} by shrinking a \textbf{loss}. Work with your
partner, show each step, and \emph{explain} what the ReLU did.
\end{headlinebox}

\vspace{0.10in}

\begin{tcolorbox}[colback=white, colframe=black!40, title=\textbf{Tier R --- Remediate},
    fonttitle=\bfseries, arc=1.5mm, left=3mm, right=3mm, top=2mm, bottom=2mm, breakable]
Warm up the two moves: the ReLU, then one layer.
\begin{enumerate}[leftmargin=1.7em, itemsep=8pt]
  \item Evaluate the ReLU (keep positives, zero out negatives):
        \[ \relu(4)={\color{keyred}\mathbf{4}},\ \ \relu(-3)={\color{keyred}\mathbf{0}},\ \ \relu(0)={\color{keyred}\mathbf{0}},\ \ \relu(7)={\color{keyred}\mathbf{7}}. \]
  \item \textbf{One layer.} For $A=\begin{bsmallmatrix} 2 & -1\\ -1 & 2 \end{bsmallmatrix}$ and
        $\vv=\begin{bsmallmatrix} 1\\ 3 \end{bsmallmatrix}$ (take $\bb=\zero$), compute the linear step
        $A\vv=\ans{$\begin{bsmallmatrix} -1\\ 5 \end{bsmallmatrix}$}$, then apply ReLU to each entry:
        $\relu(A\vv)=\ans{$\begin{bsmallmatrix} 0\\ 5 \end{bsmallmatrix}$}$. Which entry did
        the ReLU change? \ansline{The first entry: the ReLU turned $-1$ into $0$; the positive $5$ is unchanged.}
\end{enumerate}
\end{tcolorbox}

\begin{tcolorbox}[colback=white, colframe=black!40, title=\textbf{Tier A --- Approaching Proficiency},
    fonttitle=\bfseries, arc=1.5mm, left=3mm, right=3mm, top=2mm, bottom=2mm, breakable]
Now a layer with a bias, a bending function, and a loss.
\begin{enumerate}[leftmargin=1.7em, itemsep=9pt]
  \item \textbf{A full layer.} For $A=\begin{bsmallmatrix} 1 & 2\\ 1 & -1 \end{bsmallmatrix}$,
        $\bb=\begin{bsmallmatrix} -5\\ 0 \end{bsmallmatrix}$, and $\vv=\begin{bsmallmatrix} 2\\ 1 \end{bsmallmatrix}$,
        compute $A\vv$, then $A\vv+\bb$, then $\relu(A\vv+\bb)$.
        \ansline{$A\vv=\begin{bsmallmatrix} 4\\ 1 \end{bsmallmatrix}$; $A\vv+\bb=\begin{bsmallmatrix} -1\\ 1 \end{bsmallmatrix}$; $\relu(A\vv+\bb)=\begin{bsmallmatrix} 0\\ 1 \end{bsmallmatrix}$ (the bias pushed the first entry negative, then ReLU zeroed it).}
  \item \textbf{Bending with ReLUs.} Let $F(x)=\relu(x)+\relu(x-2)$. Evaluate $F(1)$ and $F(3)$. The graph
        has slope $0$ before $x=0$, slope $1$ between the two folds, and slope {\color{keyred}\textbf{2}}
        after $x=2$. Where are the two folds?
        \ansline{$F(1)=\relu(1)+\relu(-1)=1+0=1$; $F(3)=\relu(3)+\relu(1)=3+1=4$. The folds are at $x=0$ and $x=2$ (where each ReLU switches on).}
  \item \textbf{Comparing losses.} The loss is the sum of squared errors. Targets are $4,3,4$.
        Network~1 predicts $5,3,6$; Network~2 predicts $4,3,5$. Compute both losses. Which network fits
        better --- and would learning move \emph{toward} or \emph{away} from those weights?
        \ansline{Loss$_1=(5{-}4)^2+(3{-}3)^2+(6{-}4)^2=1+0+4=5$. Loss$_2=(4{-}4)^2+(3{-}3)^2+(5{-}4)^2=0+0+1=1$. Network~2 fits better (loss $1<5$); learning moves the weights \emph{toward} the smaller loss.}
\end{enumerate}
\end{tcolorbox}

\begin{tcolorbox}[colback=white, colframe=black!40, title=\textbf{Tier E --- Extension},
    fonttitle=\bfseries, arc=1.5mm, left=3mm, right=3mm, top=2mm, bottom=2mm, breakable]
\textbf{Why the ReLU is not optional.} A network stacks layers. What if we \emph{dropped} the ReLU and
just multiplied by two matrices in a row? Use $A_1=\begin{bsmallmatrix} 1 & 1\\ 0 & 1 \end{bsmallmatrix}$
and $A_2=\begin{bsmallmatrix} 2 & 0\\ 1 & 1 \end{bsmallmatrix}$.
\begin{enumerate}[leftmargin=1.7em, itemsep=9pt]
  \item Compute the single matrix $A_2A_1$.
        \ansline{$A_2A_1=\begin{bsmallmatrix} 2 & 2\\ 1 & 2 \end{bsmallmatrix}$.}
  \item For $\vv=\begin{bsmallmatrix} 1\\ 1 \end{bsmallmatrix}$, check that $A_2(A_1\vv)$ and $(A_2A_1)\vv$
        give the \emph{same} answer.
        \ansline{$A_1\vv=\begin{bsmallmatrix} 2\\ 1 \end{bsmallmatrix}$, then $A_2\begin{bsmallmatrix} 2\\ 1 \end{bsmallmatrix}=\begin{bsmallmatrix} 4\\ 3 \end{bsmallmatrix}$; and $(A_2A_1)\vv=\begin{bsmallmatrix} 4\\ 3 \end{bsmallmatrix}$. Same.}
  \item So two linear steps with no ReLU collapse into \emph{one} matrix --- a single straight-line map
        that cannot bend. In one sentence, what job does the ReLU do that makes a network more than one
        matrix?
        \ansline{The ReLU's fold makes the layer \emph{nonlinear}, so stacked layers no longer multiply out to a single matrix --- that bending is what lets a network fit curved data.}
\end{enumerate}
\end{tcolorbox}

\begin{teachernote}
Tier~R drills the two moves in isolation: ReLU arithmetic and one linear step $A\vv$ followed by ReLU
(item~2 makes the ReLU visibly zero a negative entry). Tier~A is the heart: item~1 adds a bias (which can
push an entry negative so ReLU acts); item~2 shows folds combining into a bending graph (Lesson~8.1
preview); item~3 uses the loss to \emph{compare} two networks and names the direction learning moves.
Tier~E is the key conceptual payoff --- without a nonlinear step, $A_2A_1$ is just one matrix (Unit~1
matrix multiplication), so the ReLU is exactly what makes a deep network more powerful than a single
linear map. If time is short, Tier~R item~2 and Tier~A items~2--3 are the must-dos. Watch for students who
apply ReLU to the vector as a whole instead of entry-by-entry, or who forget to square in the loss.
\end{teachernote}

\end{document}
